% 自编教学补充；阅读来源见本章 README。
\begin{frame}[fragile]{4A 标签切片与位置切片}
\small 学号设为索引后，loc 用学号找人，iloc 按当前位置找人。在此唯一且有序的索引上，loc 切片包括末端标签，iloc 不包括末端位置。
\medskip
\begin{lstlisting}
by_id = students.set_index("学号")
print(by_id.loc["001":"003", ["姓名", "成绩"]])
print(by_id.iloc[0:2, [0, 2]])
print(by_id.reset_index().columns.tolist())
\end{lstlisting}
\end{frame}

\begin{frame}[fragile]{4B 条件筛选与安全赋值}
\small 用 loc 一次指定行与列进行赋值；不要在 df[条件][列名] 的临时结果上修改原表。保留原始成绩，另建调整列。
\medskip
\begin{lstlisting}
students["调整成绩"] = students["成绩"]
mask = students["班级"].eq("B")
students.loc[mask, "调整成绩"] = students.loc[mask, "成绩"] + 2
selected = students.loc[students["成绩"].ge(80),
                        ["姓名", "班级", "成绩"]].copy()
print(selected)
\end{lstlisting}
\end{frame}
\begin{frame}[fragile]{课堂练习：4B 条件筛选与安全赋值}
筛选 A 班且成绩至少为 90 的学生，仅展示姓名与成绩。
\medskip
\begin{block}{检查思路}
先写出预期结果，再运行。参考答案见配套 Notebook 的折叠单元格。
\end{block}
\end{frame}

\begin{frame}[fragile]{4C 对齐：标签相同才相加}
\small Series 运算按标签匹配。缺少对应项会得到缺失值；fill\_value=0 只适合“未出现确实表示没有”的业务含义。
\medskip
\begin{lstlisting}
morning = pd.Series([3, 5], index=["茶", "咖啡"])
afternoon = pd.Series([2, 4], index=["咖啡", "果汁"])
print(morning + afternoon)
print(morning.add(afternoon, fill_value=0))
\end{lstlisting}
\end{frame}
